\section{背景与相关工作}
\label{sec:background}

\paragraph{定位器脆性与 benchmark。}
ReproBreak~\cite{mouraReproBreak2026} 是迄今最大的可复现定位器失效语料（来自 359 个带 Cypress/Playwright 测试的 Web 仓库的 449 个可复现失效）。我们将 \texttt{tryghost/koenig}~\cite{koenig} 的 Playwright 切片（93 个可复现失效，分布于 19 个 commit）用作主要基底。

\paragraph{自愈工具。}
Healenium~\cite{healenium} 在运行时通过 DOM 相似度修复 Selenium 选择器，是最流行的实务界 baseline。Joseph~\cite{joseph2026zerocost} 通过 DOM 可访问性树确定性地重新抽取选择器，在单个电商 demo 上报告了 100\% 通过率；这是确定性静态修复轴上最强的已发表 baseline，也是与我们 L1 基底最直接的比较对象。我们尚未注意到有任何同行评议的 LLM 定位器修复工具在真实 Playwright benchmark 上报告了误修复率；与之最接近的当代工作关注的是修复成功率或解释一致性指标，二者都不直接约束我们探针所针对的失败模式。

\paragraph{缺失的部分。}
上述工作均未在真实 benchmark 上隔离\emph{隐性误修复}—— 即修复器给出语法通过但错指元素的选择器、从而遮蔽真实 UI 变更的比率。Joseph 在一个右侧答案按构造一定可达的 demo 上报告了通过率；ReproBreak 提供了数据基底但未在其上评估任何修复器。\S\ref{sec:findings} 通过一个对抗性构造的探针填补这一空缺。

\paragraph{React 相关细节。}
\texttt{koenig-lexical}~\cite{lexicalRichText, koenig} 是一个基于 React + Lexical 的现代富文本编辑器，其 Playwright 套件在很大程度上依赖自定义 \texttt{data-kg-*} 属性（约占其失效的一半），而非 \texttt{data-testid}。\texttt{payloadcms/payload}~\cite{payloadCMS} 是一个 Next.js 后台 UI，其 CSS 类名遵循 BEM~\cite{bemMethodology}，使用 \texttt{\$\{baseClass\}\_\_suffix} 模板字面量 —— 这些字面量对 React 可见，但对纯字面量的 AST 抽取器不可见。这两种锚点文化的对比是我们跨应用发现的依据（\S\ref{sec:findings}，F3）。
